
de @@M0@@, qui résulte de la connexité géométrique de @@M1@@.

De l'homomorphisme (7.12.1) on déduit grâce à 7.4 un homomorphisme
@@M2@@

Lorsqu'on suppose de plus @@M3@@ et @@M4@@ réguliers, alors (7.12.1) peut encore s'interpréter comme un homomorphisme d'anneaux
@@M5@@
qui, lorsque @@M6@@ est de plus séparé, est un homomorphisme de @@M7@@-anneaux, et définit par suite un homomorphisme
@@M8@@

\textbf{7.12.5.} Soit @@M9@@ un schéma algébrique lisse et propre sur un corps @@M10@@. On vérifie facilement qu'un élément @@M11@@ de @@M12@@, tel qu'il existe un entier @@M13@@ avec @@M14@@ algébriquement équivalent à zéro, est numériquement équivalent à zéro. Il est plausible que la réciproque soit toujours vraie, du moins si @@M15@@ est projectif. Dans cette question, on peut manifestement supposer @@M16@@ algébriquement clos, et remplacer @@M17@@ par @@M18@@ ou @@M19@@, sans modifier la conjecture. On trouve également une conjecture équivalente en travaillant avec l'anneau de Chow @@M20@@ (cf. XIV 4.1), et on trouve la conjecture plus ou moins classique suivante : si @@M21@@ est un cycle algébrique numériquement équivalent à zéro sur @@M22@@, il existe un entier @@M23@@ avec @@M24@@ algébriquement équivalent à zéro. Dans le cas où cette conjecture est correcte pour @@M25@@, il résulterait du résultat de compatibilité de l'homomorphisme (7.3.1) avec l'équivalence algébrique que cet homomorphisme de spécialisation (dans le cas où @@M26@@ est un morphisme projectif et lisse) est aussi compatible avec l'équivalence numérique, comme conjecturé dans 7.9.

